%% Section 9: Conclusion. (outline v2 sec.3.9)

Revision-generative serving turns generated outputs into speculative writes
to future model state. This paper formalized the resulting contamination
channel and introduced certificate-bearing transactions that separate
display authority from state-write authority. Under the stated trust and
binding assumptions, the protocol excludes uncertified ancestry, produces a
single fork-free committed revision history, and remains safe for any
granularity ratio. Its cost is governed by the complete authority path---not
granularity alone---and the measurements identify when coarse authority is
worthwhile. Across two architecture families and five checkpoints, RevTrace
reaches up to $6.88\times$ wall-clock speedup. We observe no output/state or
transaction-install violations in 800 Qwen revision events (480 SG and 320
DGFT) and 1{,}048 Whisper commits; these finite traces support the declared
contracts but are not a proof of zero implementation faults. Same-granularity
certification is the $\granratio=1$ reduction of the protocol. Risk signals
schedule verification but never authorize installation, and quality evidence
does not transfer across estimands or canonical paths.
On an independent single-node 8$\times$V100 contract, balanced
four-proposer/four-authority scaling sustains $147.1\times/114.5\times$ audio
real-time throughput for Qwen3-ASR 0.6B/1.7B at c64, with display P99
1.115/1.471\,s and nominal scale efficiency 91.1/100.7\%. The c128 points fail
the 2\,s SLO and are retained as the measured capacity boundary. These values
extend the evidence to high-concurrency single-node serving without being
pooled with the RTX~3090/vLLM contract.
Implementing the session-partitioned ledger across nodes and measuring its
consensus/commit cost, and extending the transaction discipline to other
feedback-driven text and multimodal revision models, are the natural next
steps.
